% generated by roam.py handbook export -- не редактировать, править заметки
% 20260918142926-наименьшее_значение_выпуклои_функции_достигается_в_точке_с_нулевым_градиентом.org (body 0d71ae19db42)

\begin{theorem}[{Наименьшее значение выпуклой функции достигается в точке с нулевым градиентом}]\label[theorem]{note:035B8935-65A9-4BB4-BDA3-60A60FF278F0}
  Пусть $E \subset \R^n$~--- \hyperref[note:316B5BFA-8B9F-4267-9086-D6593AAB5F9B]{выпуклое множество}, $f: E \rightarrow \R$~--- \hyperref[note:1CA14F2C-637F-41DE-9E97-72075D19E97F]{выпуклая} и дифференцируемая;
  возьмем  $x^{*} \in E$
  Тогда $\nabla f(x^{*}) = 0 \Leftrightarrow x^{*}$~--- точка минимума ($\nabla f$~--- \hyperref[note:1756907B-1FF2-4F62-8553-759B8509281C]{градиент})
\end{theorem}

\begin{remark}
  Для $\Leftarrow$ нужно, чтобы $x^{*} \in \operatorname{Int} E$
\end{remark}

\begin{proof}
  \proofpart{$\Rightarrow$}
  Возьмем неравенство из \hyperref[note:18BD8497-CEBD-4962-8144-474B489606C3]{дифференциального критерия выпуклости}
  \[
  f(y) \ge f(x) + \left\langle \nabla f(x); y-x \right\rangle\ \forall x, y \in E
  \]
  Если подставить $x^{*}$, вместо $x$, получим
  \[
  \forall y \in E\ f(y) \ge f(x^{*}) + \left\langle \nabla f(x^{*}); y - x^{*} \right\rangle
  \Rightarrow f(y) \ge f(x^{*}) + \left\langle 0; y-x^{*} \right\rangle
  \]
  \[
  \Rightarrow f(y) \ge f(x^{*})\ \forall y \in E
  \]

  Причем оно достигается, так как $x^{*} \in E$.

  \proofpart{$\Leftarrow$}
  Давайте возьмем произвольное направление $h\in \R^n$,
  Так как $x^{*}$~--- внутренняя точка $E$,  найдется окрестность этой точки~--- $U$, такая, что $U \subset E$.
  А из этого следует, что $\exists \varepsilon >0:\ x^{*} + th \in U,\ |t| < \varepsilon$
  Рассмотрим $g(t) = f(x^{*} + th)$, она определена на $(-\varepsilon; \varepsilon)$.
  Так как $x^{*}$~--- точка минимума $f$, то $0$~--- точка минимума $g$, так как $g(0) = f(x^{*})$
  $g$~--- дифференцируема. Тогда по \hyperref[note:1E03F05F-22CC-4D55-BCF1-EBB5774765F0]{теореме Ферма} $g^{\prime}(0) = 0$.

  Что представляет из себя $g^{\prime}(t)$:
  \[
  g^{\prime}(t) = \frac{\partial }{\partial t} f (x^{*} + th) = f^{\prime}(x^{*} + th) h = (\nabla f( x^{*} + th))^T h = \left\langle
  \nabla f(x^{*} +th); h
  \right\rangle
  \]

  Но мы доказали, что для произвольного  $h$, $g^{\prime}(0) = \left\langle \nabla f (x^{*}); h \right\rangle = 0$
  Получается, если взять $h = \nabla f(x^{*})$ и подставить в производную, то
  \[
  \left\langle \nabla f(x^{*}) ; \nabla f(x^{*}) \right\rangle = \| \nabla f (x^{*}) \| ^2 = 0 \Rightarrow \nabla f(x^{*}) = \overline{0}
  \]
\end{proof}
